\emph{Preliminaries.} From Theorem \claimref{synthesis.chsh-bound-for-two-subsystems-of-a-clifford-algebra-by-a-rational-certificate} (the base), whose algebra is that of Theorems \claimref{synthesis.signed-directions-and-the-integer-clifford-algebra} and \claimref{synthesis.division-gives-the-rationals}: $\dagger$ is an anti-automorphism, states exist, and for a state and self-adjoint $h$, $\omega(h)^{2} \le \omega(h^{2})$ and $\omega(h^{2}) = \omega(hh^{\dagger}) \ge 0$; a functional $\langle\rho\,\cdot\,\rangle_{0}/\langle\rho\rangle_{0}$ with $\rho = \sum_{r}y_{r}y_{r}^{\dagger} \ne 0$ is a state.

\emph{(i).} $A_{+}A_{-} = A_{1}^{2} - A_{1}A_{2} + A_{2}A_{1} - A_{2}^{2} = -[A_{1},A_{2}]$ and $A_{-}A_{+} = [A_{1},A_{2}]$; $A_{\pm}^{2} = 2 \pm (A_{1}A_{2} + A_{2}A_{1})$. Regrouping, $C_{AB} = (A_{1} + A_{2})B_{1} + (A_{1} - A_{2})B_{2}$, and likewise for $D$. Put $U = A_{+}B_{1}$, $V = A_{-}D_{2}$. Since $B_{1}$ and $D_{2}$ commute with $A_{1}, A_{2}$ and with each other, $U^{\dagger} = B_{1}A_{+} = U$, $V^{\dagger} = V$, $U^{2} = A_{+}^{2}$, $V^{2} = A_{-}^{2}$, and $UV + VU = (A_{+}A_{-} + A_{-}A_{+})B_{1}D_{2} = 0$. So $(\lambda U + \mu V)^{2} = \lambda^{2}A_{+}^{2} + \mu^{2}A_{-}^{2}$. The same holds for $U' = A_{-}B_{2}$, $V' = A_{+}D_{1}$ with $A_{+}$ and $A_{-}$ exchanged.

\emph{(ii).} Let $a = \omega(A_{+}^{2})$; then $a \ge 0$ and $4 - a = \omega(A_{-}^{2}) \ge 0$. Write $u = \omega(U)$, $v = \omega(V)$, $u' = \omega(U')$, $v' = \omega(V')$, so that $\omega(C_{AB}) = u + u'$ and $\omega(C_{AD}) = v' + v$. By (i) and the preliminaries, for all scalars $\lambda, \mu$,
\[
(\lambda u + \mu v)^{2} \le \lambda^{2}a + \mu^{2}(4 - a), \qquad (\lambda u' + \mu v')^{2} \le \lambda^{2}(4 - a) + \mu^{2}a . \tag{$*$}
\]
\emph{Case $0 < a < 4$.} Put $\tau = u^{2}/a + v^{2}/(4 - a)$ and take $\lambda = u/a$, $\mu = v/(4 - a)$ in the first inequality: $\tau^{2} \le \tau$, so $\tau \le 1$. Likewise $\tau' = u'^{2}/(4 - a) + v'^{2}/a \le 1$. The identity
\[
4\Bigl(\frac{u^{2}}{a} + \frac{u'^{2}}{4 - a}\Bigr) - (u + u')^{2} = \frac{\bigl((4 - a)u - au'\bigr)^{2}}{a(4 - a)} \ge 0
\]
and the same identity for $v'$, $v$ (with $a$ under $v'$) give $(u + u')^{2} + (v + v')^{2} \le 4(\tau + \tau') \le 8$. \emph{Case $a = 0$.} From ($*$): $\mu = 0$ gives $u = 0$, and $\lambda = 0$ gives $v^{2} \le 4$; in the second, $\lambda = 0$ gives $v' = 0$ and $\mu = 0$ gives $u'^{2} \le 4$. So the sum of squares is $u'^{2} + v^{2} \le 8$. \emph{Case $a = 4$} is the same with the roles exchanged. If both values exceeded $2$ in absolute value the sum of squares would exceed $8$. Only positivity and normalisation of $\omega$ were used.

\emph{(iii).} The first sentence is clause (iii) of the base. In the commutative case each observable is $\sum_{i}\epsilon_{i}q_{i}$ with $\epsilon_{i}^{2} = 1$ (its square is $1 = \sum q_{i}$), so $C_{AB} = \sum_{i}\gamma_{i}q_{i}$ with $\gamma_{i} = \alpha_{1}(\beta_{1} + \beta_{2}) + \alpha_{2}(\beta_{1} - \beta_{2})$ for signs $\alpha, \beta$; one of $\beta_{1} \pm \beta_{2}$ is $0$ and the other $\pm 2$, so $\gamma_{i} = \pm 2$. $\omega(q_{i}) = \omega(q_{i}q_{i}^{\dagger}) \ge 0$ and $\sum_{i}\omega(q_{i}) = 1$. For commuting observables $A$, $B$ the product $AB$ is an observable, so $\omega(AB)^{2} \le \omega(1) = 1$, and $|\omega(C_{AB})| \le 4$ by the triangle inequality. The last sentence cites clauses (iii), (v) and (vi) of the base.

\emph{(iv).} The adjoint fixes $e_{j}$ and negates $f_{j}$; $p_{j}$ is self-adjoint of square $1$ and anticommutes with $e_{j}$, $f_{j}$, and commutes with the generators of the other modes. The twists are those of Theorem \claimref{tower.commuting-copies} with the plane of mode $1$ and with the planes of modes $1, 2$; the lines needed are these. $e_{2}p_{1}$ has square $e_{2}^{2}p_{1}^{2} = 1$ ($p_{1}$ commutes with $e_{2}$) and commutes with $e_{1}$, $f_{1}$ (each anticommutes with $e_{2}$ and with $p_{1}$), as in clause (vii) of the base. $p_{1}p_{2}$ has square $1$, commutes with $e_{3}$, and anticommutes with each of $e_{1}, f_{1}, e_{2}, f_{2}$; so $e_{3}p_{1}p_{2}$ has square $1$ and commutes with each of these four generators (two sign changes), hence with the algebra of modes $1, 2$, which contains $e_{1}, p_{1}, e_{2}p_{1}, p_{2}$; $p_{3}$ is even in mode $3$ and commutes with modes $1, 2$. Both twisted elements are self-adjoint (a product of commuting self-adjoint factors). $e_{3}p_{1}$ anticommutes with $e_{2}p_{1}$, since $e_{3}$ and $e_{2}$ anticommute. The total parity $p_{1}p_{2}p_{3}$ anticommutes with $e_{1}$, $e_{2}p_{1}$ and $e_{3}p_{1}p_{2}$. The three factors of $\rho$ are $1 + g$ for commuting observables $g$ (the check verifies this and the values), so $\rho$ is a positive multiple of a self-adjoint idempotent; with $B_{1} = B_{2}$, $C_{AB} = 2A_{1}B_{1}$ and $\omega_{\rho}(A_{1}B_{1}) = 1$, likewise for $D$. The state of $y\rho y^{\dagger}$ is a state by the preliminaries; its two values are computed by the check. $\rho_{0}$ is a product of two commuting self-adjoint idempotents, $y'^{\dagger}y' = 1 + \tfrac14 = \tfrac54$ ($I_{B}^{\dagger} = -I_{B}$, $I_{B}^{2} = -1$), so $\rho_{AB} = y'\rho_{0}y'^{-1}$ is a self-adjoint idempotent; its value $\tfrac{14}{5}$ is computed by the check. $p_{3}$ commutes with modes $1, 2$, so $\rho_{AB}\cdot\tfrac12(1 + p_{3})$ is a self-adjoint idempotent with the same value of $C_{AB}$. The bound $\tfrac{4}{25}$ is (ii). For this instance and a direction $a$ of mode $1$, $d = x\,e_{3}p_{1}p_{2} + z\,p_{3}$ of party $D$: the term with $e_{3}$ has no scalar part, because no blade of $\rho_{AB}\cdot\tfrac12(1 + p_{3})$ or of $a$ contains $e_{3}$ without $f_{3}$; the term with $p_{3}$ has scalar part $\tfrac{z}{2}\langle\rho_{AB}\,a\rangle_{0}$, and $\langle\rho_{AB}\,e_{1}\rangle_{0} = \langle\rho_{AB}\,p_{1}\rangle_{0} = 0$ (computed by the check; the blades of $\rho_{AB}$ other than $1$ are $e_{1}f_{2}$ and products $A_{i}B_{j}$). So all values of $C_{AD}$ are $0$ for this instance, which the check asserts; it is a trivial sample of the bound $\tfrac{4}{25}$.

\emph{(v).} The first statement is (ii). \emph{The pair state.} $\rho_{c,s}$ is self-adjoint ($\dagger$ exchanges $v_{j}$ and $w_{j}$, so $(w_{1}w_{2})^{\dagger} = v_{2}v_{1}$). With $m = cs$, $\rho_{c,s}^{2} - \rho_{c,s}$ is a polynomial of degree at most $2$ in $(r, m)$ with coefficients in the algebra, and the amplitudes lie on the conic $r^{2} + 4m^{2} = 1$; a polynomial of degree at most $2$ that vanishes at five points of an irreducible conic vanishes on it, and the check verifies $\rho_{c,s}^{2} = \rho_{c,s}$ at five rational points of the conic. $\langle\rho_{c,s}\rangle_{0} = \tfrac14$ since the other terms are not scalars. So $\omega(x) = 4\langle\rho_{c,s}x\rangle_{0}$ is a state. \emph{Correlations.} The expression is $\rho = \rho^{(0)} + r\rho^{(1)} + cs\,\rho^{(2)}$ with fixed elements $\rho^{(i)}$, so each correlation is $\alpha + \beta r + \gamma\,cs$ with rational constants, and the check evaluates all four at three weight pairs whose vectors $(1, r, cs)$ are linearly independent, which determines the constants: $2cs$, $1$, $0$, $0$ (these are the values $2cs$ and $1$ are those of the amplitudes $c$ on $\ket{00}$ and $s$ on $\ket{11}$: $e_{1}\,e_{2}p_{1}$ exchanges the two terms and $p_{1}p_{2}$ fixes both). With $B_{\pm} = zp_{2} \pm x\,e_{2}p_{1}$: $C = p_{1}(B_{+} + B_{-}) + e_{1}(B_{+} - B_{-}) = 2z\,p_{1}p_{2} + 2x\,e_{1}\,e_{2}p_{1}$, of value $2z + 4cs\,x$; by the Cauchy inequality $(z + 2cs\,x)^{2} \le (z^{2} + x^{2})(1 + 4c^{2}s^{2})$, and $4c^{2}s^{2} = 1 - r^{2}$. The numbers at the stated point are arithmetic.
